\section{Discussion}\label{discuss}

\subsection{What the probability diagnostic adds}

The main contribution is not a new HVAC method. It is a diagnostic interface for reading a calibrated TSV probability distribution. Expected TSV, defined as the probability-weighted mean $\mu_{\mathrm{TSV}}=\sum_k k p_k$, remains useful: in the full panel it was strongly associated with $p_{\mathrm{tail}}$, and the diagnostic did not produce near-neutral high-tail states under $|\mu_{\mathrm{TSV}}|<0.15$. This is an important constraint on the interpretation. The probability representation does not overturn the scalar mean.

The added value is that the distribution makes tail probability explicit. A scalar mean reports central tendency; $p_{\mathrm{tail}}$ reports the probability mass in the selected cool/cold and warm/hot categories. Similar mean states can still have different tail probabilities, and the threshold status of $p_{\mathrm{tail}}$ is directly available only when the distribution is retained. In this fixed-reference future-weather panel, the visible tail was also strongly warm-skewed: cold-tail probability was small, and high-tail records were warm-dominant. That pattern is conditional on the selected cities, weather files, fixed thermostat schedule, and static comfort model, but it is exactly the kind of directional diagnostic that the probability object makes explicit.

The panel summary in Table~\ref{tab:panel_tail_summary} also matters because it separates broad stress-test gradients from local case differences. High-tail exceedance was similar in the reference and near-2030s slices, then increased in the mid-2050s and late-2080s; it was also higher under SSP5-8.5 than SSP2-4.5. The typical-year selector followed the same time-slice ordering, while the hot and heatwave-extreme annual files were close in aggregate exceedance. The city spread was larger than either scenario contrast or annual-severity contrast. This supports a cautious reading: the fixed archetype shows more high-tail states in the later selected weather, but the diagnostic should not be reduced to a single climate-scenario ranking. Local climate, selected annual file, and building-zone response still shape the observed tail profile.

The weather-to-zone audit supplies a physical pathway for that descriptive gradient. All 36 matched reference-to-late role comparisons moved toward warmer outdoor and simulated operative conditions and higher mean tail probability, and the direction also held in all 12 city--scenario cells after duplicate weather roles were collapsed. The cooling system responded concurrently, but the available traces do not establish equipment saturation. Likewise, the south-perimeter concentration is consistent with orientation, envelope, and zoning interaction in this prototype, but horizontal GHI cannot identify facade-incident solar causation. These safeguards are important because the fitted probability model also conditions on humidity and running-mean outdoor context: whole-year rank associations describe seasonal co-movement, not a causal temperature coefficient.

\subsection{Interpretation of the discomfort-relevant TSV tail}

The \(|\mathrm{TSV}|\ge2\) category boundary is not an assertion that every warm or cool vote is unacceptable. It is a convention-anchored screening construct. On the seven-point scale, \(+1\) is ``slightly warm'' and remains in the central band used here; \(+2\) is ``warm.'' The PMV--PPD lineage associates the \(\pm2\) and \(\pm3\) sensation categories with dissatisfaction, which makes the selected tail relevant to comfort screening \citep{Fanger1970PMV,ASHRAE2023}. That historical grouping anchors the boundary but does not turn a sensation vote into a directly observed satisfaction, preference, or acceptability response.

This distinction is especially important under adaptive comfort theory. Occupants with adaptive opportunity may accept or prefer a non-neutral sensation, including a warm state. Accordingly, the present study does not interpret \(p_{\mathrm{tail}}\) as an acceptability probability or treat the 0.20 reporting screen as a standards threshold. Formal adaptive-comfort compliance is also not the applicable normative path for the mechanically conditioned fixed-schedule prototype. The adaptive evidence is used here to bound interpretation, not to erase the diagnostic relevance of probability outside the central \(-1,0,+1\) sensation band.

\subsection{PMV is a comparator, not a strawman}

The PMV results also support a restrained interpretation. PMV was correlated with $p_{\mathrm{tail}}$, and the standard $|\mathrm{PMV}|=0.5$ band was conservative in this fixed-reference panel. It almost never missed high-tail states. Its main difference from $p_{\mathrm{tail}}$ was over-flagging: many states outside the PMV neutral band did not cross the selected TSV tail-probability cutoff. The comparison is also geometrically informative. The learned expected-TSV scalar stays within a narrower empirical range and tracks the calibrated tail probability closely, while PMV spans a wider continuous comfort-index range and includes a lower-tail-probability branch at elevated $|\mathrm{PMV}|$. This tighter mapping is unsurprising because $\mu_{\mathrm{TSV}}$ and $p_{\mathrm{tail}}$ are two summaries of the same calibrated TSV probability vector; it should not be interpreted as a claim that the relationship is linear. That branch means that some states can look substantially displaced on the PMV scale while the calibrated TSV distribution still assigns most probability to neutral or mild bins. This is expected: $\mu_{\mathrm{TSV}}$ is a probability-weighted mean of the calibrated seven-class TSV predictor used in this study, whereas PMV is Fanger's continuous index built from thermal-balance terms and empirical chamber-vote evidence. The two scalars therefore stretch the same simulated thermal states differently.

Table~\ref{tab:pmv_scalar_comparison} makes the threshold consequence explicit. An offline optimized $|\mathrm{PMV}|$ threshold can be made closer to the $p_{\mathrm{tail}}$ screen, but it still disagrees more often than the optimized expected-TSV threshold. The conventional $|\mathrm{PMV}|=0.5$ band has a different purpose: it is protective as a comfort band, not tuned to reproduce this paper's tail-probability screen. The high false-positive rate under that band is therefore not a defect by itself; it shows that a conventional scalar comfort band and the selected TSV-tail probability flag are not interchangeable diagnostics.

This means the diagnostic should not be framed as PMV failure. PMV remains a valuable standard comfort index and benchmark. The primary TSV model also uses PMV as one covariate, so the PMV panel is not an independent model contest. The no-PMV robustness check is therefore important, but it has a narrow meaning: the same synchronized PMV scalar was compared against two different TSV probability outputs, one from the primary model and one from a separately trained model without PMV. Removing PMV modestly loosened the PMV--$p_{\mathrm{tail}}$ relationship, which is consistent with PMV carrying useful information for reported TSV prediction, but it left the main aggregation result nearly unchanged. The probability diagnostic answers a different question: how much of the calibrated empirical TSV distribution lies in the selected TSV tail? For building studies that need a compact scalar, PMV or expected TSV may be enough. For studies that need category-probability screening, the probability vector carries information that scalar summaries do not directly expose.

\subsection{Zone aggregation is a major source of information loss}

The zone results show that aggregation can matter more than the choice between two scalar summaries. The mean-zone screen was 13.15\%; the area-weighted mean-probability screen and area-weighted zone-time screened share were 12.13\% and 12.87\%, respectively. Yet localized screening found at least one high-tail zone in 23.5\% of occupied records that the mean-zone summary missed. This is not a minor reporting difference; it changes the screened profile of the same fixed simulated building. The threshold sweep strengthens this point. Changing the high-tail screen changes the absolute exceedance percentage, as expected, but it does not remove the gap between mean-zone, percentile-zone, and localized screening summaries. The same ordering persisted with a nominal seven-class predictor and with the endpoint-only $\{-3,+3\}$ tail, although the absolute rates were model-, threshold-, and event-definition-conditional.

The zone pattern should not be summarized as a simple ``top floor is worst'' or ``bottom floor is worst'' result. In this prototype, P1 is the south perimeter zone and is the dominant local exceedance zone; bottom P1 was higher than middle and top P1 in most cases. At the same time, area-weighted floor exceedance was usually highest on the middle floor because the large core zone and the other perimeter orientations contributed more to the floor-level average. The any-zone floor metric answers a different question: whether at least one zone on that floor crosses the high-tail screen. It therefore exposes small-zone perimeter states rather than suppressing them. This is why bottom-floor any-zone exceedance can be high even when the middle floor has the higher area-weighted mean. The conditioned top zones are also not direct roof-surface rooms; they sit below a top-floor plenum and roof assembly. The useful interpretation is therefore more specific: orientation, floor construction, zone size, aggregation rule, and local climate interact, and the probability diagnostic makes those interactions visible.

Figure~\ref{fig:zone_floor_position} also shows that floor differences are differences in degree, not separate thermal regimes. The bottom, middle, and top floor summaries are broadly comparable in magnitude, with the middle floor slightly elevated under area weighting. The tail-component bars show the same directional structure on all three floors: warm-tail probability dominates the TSV-tail signal, while cold-tail probability remains small. This supports the paper's diagnostic framing. The main finding is not that one floor is categorically unsafe, but that the conclusion changes depending on whether the analyst reports a mean, area-weighted, percentile, or any-zone view.

The result should still be read carefully. EnergyPlus zone states are well-mixed summaries, so the diagnostic does not prove that occupants at specific desks experienced radiant asymmetry, stratification, or local drafts. Any-zone exceedance is therefore a modeled zone-state screen, not a measured occupant-exposure rate. It shows that even at the modeled zone level, averaging can hide perimeter/core differences that are relevant to a probability-state diagnostic. A future implementation should report mean, percentile, and max-zone summaries together, or keep the probability object at the zone level when the building system can act by zone.

\subsection{Metabolic inputs change the TSV-tail profile}

The metabolic-input diagnostic turns a common modeling convention into a measurable sensitivity. Holding the environment fixed while changing the metabolic input produced threshold flips in 19.0\% of mean-zone records and 38.7\% of max-zone records. This does not mean that EnergyPlus loads would be unchanged in a real population; the diagnostic deliberately did not rerun internal heat gains. It isolates the representation question: how sensitive is the predicted TSV-tail probability to the representative metabolic assumption?

The answer is that the sensitivity is large enough to affect the screened TSV-tail profile. This is especially relevant because metabolic rate is often treated as a fixed office default. In the warm-dominant fixed-reference panel, lower metabolic-load assumptions generally lowered predicted tail exceedance because the same environmental state is interpreted with less occupant heat production. That is the model-output result. The demographic interpretation is more cautious. The metabolic-rate correction used to define the profiles shows that lower W/person values may correspond to population groups such as older or smaller-bodied occupants \citep{Guo2025Correcting120W}. The result should not be read as evidence that those occupants are safer or experience less actual discomfort. A recent review of elderly thermal characteristics emphasizes that older adults' reported thermal sensation and discomfort responses can differ from those of younger adults, including altered sensitivity and tolerance patterns \citep{Xu2025ElderlyThermalReview}. The present diagnostic estimates the probability of reported warm/hot or cool/cold TSV categories under a trained model; it does not measure latent discomfort, physiological heat strain, or health vulnerability. The non-monotone response between some adjacent profiles also cautions against overinterpreting the model as a smooth physiological curve. Plausible metabolic assumptions expose a spread in predicted reported-TSV tail probability that a single mean-occupant convention hides.

\subsection{Downstream use without overclaiming}

The probability object could support future operational designs, but only through a separately validated decision layer. Expected TSV or PMV could retain direction and a compact mean-state reference, while warm- and cold-tail probabilities could provide asymmetric diagnostic inputs. A candidate policy would then need to define a probability trigger or explicit discomfort cost, select an actuator response, and enforce dwell time, equipment capacity, setpoint feasibility, recalibration, and fail-safe requirements. A stochastic or chance-constrained controller could use the same quantities, but its objective and constraints would require independent validation. Before any control study, an operator can use the present diagnostic more modestly to screen weather files, zones, or occupant-assumption scenarios and identify where a scalar summary is insufficient.

Those are downstream uses, not results of this paper. The fixed-reference simulations here do not test energy savings, comfort improvement, setpoint stability, equipment wear, demand response, or field deployment. The 15-minute records are diagnostic probability records. They should not be interpreted as equipment commands.

\section{Limitations and Future Work}\label{sec:limitations}

\paragraph{Static comfort calibration}
The TSV probability model is calibrated from contemporary field observations pooled from the ASHRAE Global Thermal Comfort Database II and the Chinese Thermal Comfort Dataset and is held fixed across all future-weather files. The study does not model future changes in expectations, clothing behavior, adaptive opportunity, physiology, culture, or thermal history. The results are therefore conditional stress-test diagnostics, not forecasts of how occupants will perceive buildings in the 2050s or 2080s.

\paragraph{Dataset transport}
Contributor-grouped and overlap-excluded reciprocal ASHRAE--China source-label holdout checks both showed weaker tail-event performance than the pooled record-level split. In the reciprocal check, tail AUROC was 0.553--0.592, tail ECE was 0.096--0.124, and the prevalence gap changed direction by target source. These checks bound the transportability of the study-specific fitting pipeline; they do not constitute validation of the final pooled predictor against a wholly unused third corpus. Absolute $p_{\mathrm{tail}}$ values should therefore be interpreted as conditional on the training sources, feature specification, and calibration procedure.

\paragraph{Fixed building testbed}
The Standard-2019 Denver DOE Medium Office prototype and its Denver design-day sizing basis are held fixed across all cases, while each EPW supplies its own site location and weather. This isolates weather and representation effects for one fixed design basis, but it does not represent city-specific code compliance, local design sizing, future envelope upgrades, adaptive glazing, shading retrofits, HVAC redesign, or future operation practice. Cooling-rate associations describe coincident system response; without explicit unmet-load or component-capacity evidence, they do not establish equipment saturation.

\paragraph{Weather-panel scope}
The 144 weather files are role-labelled annual selections from a preserved CMIP-derived panel and represent 119 unique source-year trajectories because independently applied severity selectors can converge on the same year. They support structured comparisons across city, scenario, time slice, and selection role, but they are not independent climate realizations or a probabilistic climate ensemble. The 2025--2029 reference window contains five candidate years, whereas the three later windows contain ten each. This unequal opportunity can affect maximum-based hot and heatwave-extreme selection, so time-slice differences are treated as descriptive stress-test contrasts rather than causal estimates of warming alone. The paper does not estimate event frequencies or regional climate uncertainty. A unique-source-year sensitivity is reported alongside the prespecified role-weighted panel. The archive contains a matching weather-generator implementation and reproduces every selector and source-to-EPW transform; however, it does not preserve the exact six-city batch configuration or its raw CMIP6 and observational inputs. Its QA therefore establishes algorithmic lineage, exact downstream integrity, and broad plausibility, not byte-level end-to-end regeneration or external validation of the exact forecast files. Moreover, only temperature, pressure, wind speed, and GHI have documented daily-delta transformations; RH, specific humidity, and direct-normal irradiance are scenario-invariant in the preserved paired files. Scenario differences cannot therefore be decomposed into independently projected humidity or direct-beam effects.

\paragraph{Conditional uncertainty}
The panel bootstrap resamples the existing city--scenario--time-slice design blocks, and the held-out bootstrap resamples contributors. Their percentile ranges quantify sensitivity to the composition of this finite test design. They do not propagate predictor refitting, climate-model ensembles, weather-generation alternatives, building-model uncertainty, occupant adaptation, or transport to a wholly unused corpus. The reported ranges are therefore conditional robustness intervals rather than confidence bounds for future occupants, buildings, or regional climate.

\paragraph{EnergyPlus zone resolution}
EnergyPlus provides well-mixed zone-level states. The zone-aggregation diagnostic therefore quantifies information loss between modeled zones. It does not resolve sub-zone radiant asymmetry, vertical stratification, local air speed, solar patches, or occupant-location-specific exposure. Reported any-zone exceedance rates should therefore be read as modeled zone-state exceedances rather than as measured occupant-exposure rates.

\paragraph{TSV-tail calibration}
The outer TSV categories are data-poor relative to the neutral class, especially at the $\mathrm{TSV}=-3$ and $+3$ endpoints. Post-hoc isotonic calibration improves probability reliability, but tail estimates remain more uncertain than central categories. Endpoint-only results are therefore reported as a support-limited definition sensitivity rather than as a more valid estimate of actual discomfort. The nominal sensitivity preserved the central spatial ordering but differed from the ordinal model in absolute screened shares, especially at lower probability thresholds; this confirms that absolute values remain conditional on predictor specification. Future work should test tail-weighted training, partial proportional-odds models, mixture-of-threshold models, and calibration methods that better preserve class dependence in sparse tails.

\paragraph{Metabolic-input interpretation}
The metabolic diagnostic recomputes TSV probabilities over fixed environmental traces. It does not rerun EnergyPlus with altered occupant sensible heat gains, nor does it model correlations among metabolic rate, clothing, activity, age, sex, body size, or adaptation. The target is also reported TSV, not latent physiological strain or actual health risk. If some groups systematically report thermal sensation or discomfort differently from younger reference populations, or if TSV datasets encode reporting and adaptation biases, the diagnostic would inherit those biases. Lower predicted reported-TSV tail probability for a lower metabolic-rate profile should therefore not be equated with lower vulnerability. A full population model should represent these variables jointly rather than as independent scenario shocks and should distinguish reported sensation, comfort, preference, and physiological heat stress.

\paragraph{No closed-loop validation}
The present study intentionally excludes closed-loop operating validation. The diagnostic does not establish better HVAC performance, lower EUI, improved comfort, setpoint feasibility, equipment stability, or reduced demand. Those claims would require separate closed-loop experiments with explicit objectives, equipment constraints, dwell-time logic, and field or high-fidelity validation.

\section{Conclusions}

This paper presents a diagnostic-interface study of a calibrated seven-class TSV distribution. It is not a control study. Expected TSV describes the center of that distribution, while $p_{\mathrm{tail}}=\Pr(\mathrm{TSV}\le-2)+\Pr(\mathrm{TSV}\ge+2)$ describes the model-estimated mass in the cool/cold and warm/hot categories. The selected tail is discomfort-relevant by convention, but it is not observed dissatisfaction, acceptability, PPD, or health risk.

In a 144-role fixed-reference Medium Office panel, 13.1\% of occupied records crossed the illustrative $p_{\mathrm{tail}}\ge0.20$ screen. High-tail exceedance was similar in the reference and near-2030s slices, then increased in the 2050s and 2080s, and was higher under SSP5-8.5. PMV and expected TSV both tracked tail probability, but neither reports the full category distribution. Among the scalar and spatial summaries tested, the clearest information-loss result arose from zone aggregation: at least one high-tail zone occurred in 23.5\% of occupied records for which the mean-zone summary remained below the screen. That ordering persisted under nominal-model, threshold, unique-weather-state, and endpoint-only checks. Metabolic inputs also changed threshold status, especially under max-zone aggregation.

The defensible conclusion is narrow and useful. Calibrated TSV probabilities expose category-mass, spatial, and occupant-input profiles that scalar comfort states compress. This representation can inform subsequent measurement or operational-design studies, but it is not itself a validated operating strategy.
